\ssr{ВВЕДЕНИЕ}

Целью данной работы является выполнить сравнительный анализ рекурсивного и нерекурсивного алгоритмов.

Для достижения цели поставлены следующие задачи:
\begin{itemize}
	\itemразработать рекурсивный и нерекурсивный алгоритмы решения задачи нахождения максимального числа последовательности;
	\itemописать средства разработки и инструменты замера процессорного времени выполнения реализации алгоритмов;
	\itemреализовать разработанные алгоритмы;
	\itemвыполнить тестирование реализации алгоритмов;
	\itemвыполнить теоретическую оценку затрачиваемой реализацией каждого алгоритма памяти (для рекурсивного алгоритма — на материале анализа высоты дерева рекурсивных вызовов и оценки затрачиваемой на один вызов функции памяти);
	\itemвыполнить замеры процессорного времени выполнения реализации алгоритмов в зависимости от варьируемого входа;
	\itemоценить трудоёмкость двух алгоритмов или их реализаций в худшем случае;
	\itemсравнить результаты замеров процессорного времени и оценки трудоёмкости;
	\itemсделать выводы из сравнительного анализа реализации рекурсивного и нерекурсивного алгоритмов решения одной и той же задачи, заданной вариантом, по критериям ёмкостной эффективности (на материале теоретической оценки пиковой временной эффективности.
\end{itemize}
\clearpage
